% navimind_ieee.tex — IEEE 8-page double-column paper
% Compile: xelatex navimind_ieee.tex (or pdflatex if no CJK needed)
\documentclass[conference]{IEEEtran}

% ── Packages ──
\usepackage[utf8]{inputenc}
\usepackage[T1]{fontenc}
\usepackage{amsmath,amssymb,amsfonts}
\usepackage{algorithmic}
\usepackage[ruled,vlined,linesnumbered]{algorithm2e}
\usepackage{graphicx}
\usepackage{textcomp}
\usepackage{booktabs}
\usepackage{multirow}
\usepackage{xcolor}
\usepackage{hyperref}
\usepackage{cite}
\usepackage{balance}
\usepackage{subcaption}
\usepackage{enumitem}
\usepackage{url}

% ── TBD placeholder command ──
\newcommand{\TBD}[1]{\textcolor{red}{\textbf{[TBD: #1]}}}

% ── Math shortcuts ──
\newcommand{\vp}{\mathbf{p}}
\newcommand{\vf}{\mathbf{f}}
\newcommand{\vmu}{\boldsymbol{\mu}}
\DeclareMathOperator*{\argmin}{arg\,min}
\DeclareMathOperator*{\argmax}{arg\,max}
\DeclareMathOperator{\Var}{Var}
\DeclareMathOperator{\Entropy}{H}

\begin{document}

% ══════════════════════════════════════════════════════════════════════
% Title & Authors
% ══════════════════════════════════════════════════════════════════════
\title{NaviMind: Belief-Aware Hierarchical Scene Graphs for\\Robust Zero-Shot Object Navigation on Quadruped Robots}

\author{
\IEEEauthorblockN{Hongsen Pang and Hanzhuo Zhang}
\IEEEauthorblockA{Shanghai Inovxio Technology Co., Ltd.\\
Shanghai, China\\
\{hongsen.pang, hanzhuo.zhang\}@inovxio.com}
}

\maketitle

% ══════════════════════════════════════════════════════════════════════
% Abstract
% ══════════════════════════════════════════════════════════════════════
\begin{abstract}
We present \textbf{NaviMind}, a zero-shot object navigation system that enables quadruped robots to navigate to targets specified by natural language in unknown environments.
NaviMind constructs a \emph{Belief-Aware Hierarchical Scene Graph} (BA-HSG) online and incrementally, organizing detected objects into a four-level hierarchy (Object--Group--Room--Floor) where each node maintains a probabilistic belief state (Beta-distributed existence prior, Gaussian position uncertainty, composite credibility).
We make three key contributions:
(1)~BA-HSG with FOV-aware belief updates and graph diffusion for principled false-positive rejection;
(2)~a Fast--Slow dual-path architecture where 75\%+ of queries are resolved via sub-millisecond multi-source fusion while complex queries use hierarchical Chain-of-Thought LLM reasoning with 90\% token reduction;
and (3)~a Value-of-Information (VoI) scheduler that replaces heuristic fixed-interval triggers with Shannon-entropy-driven information gain optimization under edge-compute budget constraints.
The system runs entirely onboard a Unitree Go2 quadruped with Jetson Orin NX (16\,GB), addressing low-viewpoint ($\sim$30\,cm), visual jitter, and terrain-constrained navigation challenges.
Ablation studies confirm that BA-HSG improves false-positive rejection by \TBD{X\%} over deterministic baselines, the Fast Path achieves $>$75\% hit rate at $<$5\,ms latency, and VoI scheduling reduces unnecessary re-perception by \TBD{Y\%} compared to fixed-interval triggering.
\end{abstract}

\begin{IEEEkeywords}
Object Navigation, Belief Scene Graph, Uncertainty-Aware Planning, Vision-Language Navigation, Quadruped Robot, Zero-Shot
\end{IEEEkeywords}

% ══════════════════════════════════════════════════════════════════════
% I. Introduction
% ══════════════════════════════════════════════════════════════════════
\section{Introduction}
\label{sec:intro}

Object-goal navigation---the task of navigating to a target object specified by its semantic category or natural language description---is a fundamental capability for mobile robots operating in human environments~\cite{batra2020objectnav, yadav2022hm3d}.
Recent advances in large language models (LLMs) and vision-language models (VLMs) have opened new possibilities for zero-shot approaches that require no task-specific training~\cite{sgnav2024, esc2023, l3mvn2023}.
However, deploying these capabilities on real robots, particularly quadruped platforms in unknown environments, remains a significant open challenge.

\subsection{Motivation}

Current approaches face fundamental limitations when applied to real-world quadruped deployment:

\textbf{Offline map dependency.}
Methods like FSR-VLN~\cite{fsrvln2025} achieve 92\% success rate by constructing hierarchical scene graphs through offline scanning, making them inapplicable to unknown environments.
ConceptGraphs~\cite{conceptgraphs2024} and HOV-SG~\cite{hovsg2024} similarly assume pre-scanned environments.

\textbf{Flat scene representation.}
Systems like LOVON~\cite{lovon2025} and CoW~\cite{cow2023} use flat object-level matching without structured spatial reasoning.
When instructions involve spatial relations (``the fire extinguisher near the door'') or room-level context (``go to the office''), flat representations fail to provide hierarchical context for disambiguation.

\textbf{No dynamic adaptation.}
Most systems assume static environments~\cite{fsrvln2025, hovsg2024}.
Without mechanisms to verify and update target beliefs during navigation, false-positive detections lead to irreversible failures.

\textbf{Quadruped-specific challenges.}
Quadruped robots introduce extremely low camera viewpoint ($\sim$30\,cm vs.\ $\sim$90\,cm for wheeled robots), significant visual jitter from walking gaits, and limited field of view.
Only LOVON~\cite{lovon2025} and NaVILA~\cite{navila2025} address quadruped navigation, but neither uses structured scene graphs.

\subsection{Contributions}

We make the following contributions:

\begin{enumerate}[leftmargin=*]
\item \textbf{Belief-Aware Hierarchical Scene Graph (BA-HSG).}
A four-level scene graph (Object--Group--Room--Floor) where each node maintains a Beta-distributed existence belief, Gaussian position uncertainty, and composite credibility.
FOV-aware negative evidence and graph diffusion enable principled false-positive rejection.

\item \textbf{Adaptive Fast--Slow Dual-Path Inference.}
Multi-source confidence fusion resolves 75\%+ of queries in $<$1\,ms without LLM calls.
Complex queries trigger selective grounding ($\sim$90\% token reduction) and hierarchical Chain-of-Thought LLM reasoning.
An AdaNav entropy trigger forces slow-path escalation when candidate score entropy exceeds a threshold.

\item \textbf{VoI-Driven Reasoning Scheduling.}
The decision of when to re-perceive, invoke slow reasoning, or continue navigation is formalized as Shannon-entropy-based information gain optimization under edge-compute constraints, replacing heuristic fixed-interval triggers.
\end{enumerate}

The system is validated on a Unitree Go2 quadruped with Jetson Orin NX (16\,GB), with ablation studies quantifying each contribution.

% ══════════════════════════════════════════════════════════════════════
% II. Related Work
% ══════════════════════════════════════════════════════════════════════
\section{Related Work}
\label{sec:related}

\subsection{Scene Graph Navigation}

3D scene graphs provide structured representations encoding objects, attributes, and inter-object relationships.
Hydra~\cite{hydra2022} constructs 5-level real-time scene graphs using metric-semantic meshes.
ConceptGraphs~\cite{conceptgraphs2024} builds open-vocabulary 3D scene graphs incrementally from RGB-D, leveraging CLIP~\cite{clip2021} for semantic encoding.
HOV-SG~\cite{hovsg2024} extends this to multi-floor hierarchical graphs.

Most relevant to our work, SG-Nav~\cite{sgnav2024} constructs online hierarchical scene graphs with LLM-based Chain-of-Thought prompting, achieving state-of-the-art zero-shot performance on MP3D and HM3D ($>$10\% SR improvement).
FSR-VLN~\cite{fsrvln2025} proposes a four-level HMSG with Fast-to-Slow reasoning on a humanoid, but requires offline scanning.
DovSG~\cite{dovsg2024} addresses dynamic environments with local update mechanisms but lacks formal uncertainty modeling.

NaviMind extends SG-Nav's paradigm with three innovations: (i)~belief-aware nodes with calibrated uncertainty, (ii)~VoI-driven scheduling, and (iii)~risk-sensitive multi-hypothesis planning.

\subsection{Zero-Shot Object-Goal Navigation}

Zero-shot methods eliminate task-specific training by leveraging pre-trained foundation models.
CoW~\cite{cow2023} uses CLIP matching for zero-shot navigation.
ESC~\cite{esc2023} uses LLM captioning to guide frontier exploration.
L3MVN~\cite{l3mvn2023} prompts LLMs with object lists for frontier selection.
VLFM~\cite{vlfm2024} generates language-grounded value maps from RGB observations, achieving SOTA SPL on Boston Dynamics Spot.

Unlike flat-representation methods, NaviMind uses hierarchical scene graphs for structured context and avoids LLM calls for 75\%+ of queries via the Fast Path.

\subsection{Dual-Process Reasoning and Uncertainty}

Inner Monologue~\cite{innermonologue2023} grounds LLM planning with environment feedback.
SayCan~\cite{saycan2022} combines LLM knowledge with robotic affordances.
LERa~\cite{lera2025} introduces Learn-Execute-Replan with experience accumulation.

SG-Nav~\cite{sgnav2024} proposes graph-based re-perception with scalar credibility, but without formal uncertainty modeling.
Classical active perception uses information-theoretic criteria (entropy reduction, mutual information) for sensing action selection~\cite{howard1966voi}.
Our VoI scheduler adopts this principle, formalizing re-perception and slow-reasoning triggers as entropy-driven information gain optimization.

\subsection{Quadruped Navigation}

LOVON~\cite{lovon2025} integrates LLM planning with open-vocabulary detection on Unitree quadrupeds, using Laplacian blur filtering.
NaVILA~\cite{navila2025} proposes an end-to-end VLA model trained on human videos.
OrionNav~\cite{orionnav2024} constructs semantic scene graphs with LLM planning on quadrupeds but lacks hierarchical CoT and re-perception.
VLN-PE~\cite{vlnpe2025} reveals that quadrupeds suffer from restricted observation space and locomotion-induced visual degradation.

NaviMind is the first system combining online hierarchical scene graph construction, belief-aware updates, dual-path reasoning, and VoI scheduling on a quadruped platform.
Table~\ref{tab:comparison} summarizes the comparison.

\begin{table}[t]
\centering
\caption{Comparison with related systems.}
\label{tab:comparison}
\footnotesize
\setlength{\tabcolsep}{2pt}
\begin{tabular}{@{}lccccccc@{}}
\toprule
System & SG & Hier. & Online & Belief & Explore & Quad. & ZS \\
\midrule
SG-Nav~\cite{sgnav2024}        & \checkmark & 3-lvl & \checkmark & -- & LLM  & -- & \checkmark \\
FSR-VLN~\cite{fsrvln2025}      & \checkmark & 4-lvl & --         & -- & --   & -- & \checkmark \\
LOVON~\cite{lovon2025}          & --         & --    & --         & -- & --   & \checkmark & \checkmark \\
OrionNav~\cite{orionnav2024}    & \checkmark & flat  & \checkmark & -- & LLM  & \checkmark & \checkmark \\
VLFM~\cite{vlfm2024}            & --         & --    & --         & -- & VLM  & -- & \checkmark \\
DovSG~\cite{dovsg2024}          & \checkmark & flat  & \checkmark & -- & --   & -- & \checkmark \\
\textbf{NaviMind}                & \checkmark & 4-lvl & \checkmark & \checkmark & TSG+LLM & \checkmark & \checkmark \\
\bottomrule
\end{tabular}
\vspace{-2mm}
\end{table}

% ══════════════════════════════════════════════════════════════════════
% III. Method
% ══════════════════════════════════════════════════════════════════════
\section{Method}
\label{sec:method}

NaviMind consists of three modules (Fig.~\ref{fig:architecture}):
a \emph{Perception Module} (\S\ref{sec:bahsg}) that constructs and maintains an online BA-HSG $\mathcal{G}$;
a \emph{Planning Module} (\S\ref{sec:fastslow}) that resolves natural language instructions via Fast--Slow dual-path reasoning;
and an \emph{Exploration Module} that selects frontiers when the target is absent from $\mathcal{G}$.
All modules run onboard Jetson Orin NX with no cloud dependency except optional LLM API calls.

\begin{figure}[t]
\centering
\TBD{System architecture figure: Perception $\to$ BA-HSG $\to$ Fast--Slow Planning $\to$ VoI Scheduling $\to$ Nav2 Execution loop}
\caption{NaviMind system architecture. The perception module constructs a BA-HSG online from RGB-D and LiDAR-inertial odometry. The planning module resolves goals via Fast--Slow dual-path reasoning. VoI scheduling drives adaptive re-perception and slow-reasoning triggers.}
\label{fig:architecture}
\end{figure}

\subsection{BA-HSG Construction and FOV-Aware Update}
\label{sec:bahsg}

\subsubsection{Scene Graph Definition}

We define the hierarchical scene graph as $\mathcal{G} = (\mathcal{V}, \mathcal{E})$ with four node levels:
\begin{equation}
\mathcal{V} = \mathcal{V}^{\text{obj}} \cup \mathcal{V}^{\text{grp}} \cup \mathcal{V}^{\text{room}} \cup \mathcal{V}^{\text{floor}}
\end{equation}

\noindent\textbf{Object nodes} $v_i^{\text{obj}} = (\ell_i, \vp_i, \vf_i, s_i, n_i)$ store a semantic label $\ell_i$, 3D position $\vp_i \in \mathbb{R}^3$, CLIP feature $\vf_i \in \mathbb{R}^{512}$, best detection score $s_i$, and detection count $n_i$.
\textbf{Room nodes} aggregate semantically coherent spatial regions with inferred names (e.g., ``corridor'', ``office'').
\textbf{Group nodes} cluster related objects within rooms.
Edges include hierarchical affiliation (Object$\leftrightarrow$Group$\leftrightarrow$Room$\leftrightarrow$Floor) and spatial relations (near, on, left\_of, etc.).

\subsubsection{Node-Level Belief State}

Each object $v_i$ is augmented with three belief components:

\textbf{Existence belief} via Beta distribution:
\begin{equation}
P(\text{exists}_i) \sim \text{Beta}(\alpha_i, \beta_i), \quad \bar{P}_i = \frac{\alpha_i}{\alpha_i + \beta_i}
\end{equation}
Positive evidence (high-confidence detection with valid depth and geometric consistency) increments $\alpha_i$; negative evidence (object within camera FOV but undetected) increments $\beta_i$ with a weaker weight $\eta_{\text{neg}} = 0.5$.

\textbf{Position uncertainty} via isotropic Gaussian:
\begin{equation}
\vp_i \sim \mathcal{N}(\vmu_i, \sigma_i^2 \mathbf{I})
\end{equation}
Updated via Kalman-style fusion with depth-proportional observation noise $\sigma_{\text{obs}} = \sigma_0 + \sigma_d \cdot d$.

\textbf{Composite credibility}:
\begin{equation}
C_i = w_s \bar{P}_i + w_v A_i + w_t e^{-\Delta t_i / \tau_t} + w_g e^{-\epsilon_i^{\text{reproj}} / \kappa}
\label{eq:credibility}
\end{equation}
where $A_i = 1 - 1/n_i$ captures view diversity, $\Delta t_i$ is staleness, and $\epsilon_i^{\text{reproj}}$ is reprojection error.

\subsubsection{FOV-Aware Negative Evidence}

Unlike prior work that applies negative evidence globally~\cite{sgnav2024}, we restrict miss counting to objects within the camera's frustum.
At each frame, we project tracked object positions into the image plane and mark objects within the FOV polygon but undetected.
This prevents penalizing objects that are simply out of view, significantly reducing false negatives in the belief update.

\subsubsection{Graph Diffusion}

Credibility propagates through the hierarchy:
\textbf{Upward}---room credibility is the detection-weighted average of its objects': $C_k^{\text{room}} = \sum_{i \in \mathcal{O}_k} n_i C_i / \sum n_i$.
\textbf{Downward}---well-established rooms boost newly detected objects: $\alpha_i^{\text{new}} \mathrel{+}= \delta_{\text{room}} C_k^{\text{room}}$.
\textbf{Lateral}---spatially related objects share partial credibility evidence.

\subsection{Adaptive Fast--Slow Goal Resolution}
\label{sec:fastslow}

\subsubsection{Fast Path: Multi-Source Confidence Fusion}

The Fast Path computes a fused confidence score for each object without LLM calls:
\begin{equation}
s_{\text{fused}}(v_i, I) = w_\ell s_{\text{label}} + w_c s_{\text{clip}} + w_d s_{\text{det}} + w_r s_{\text{spatial}}
\label{eq:fusion}
\end{equation}
where $s_{\text{label}}$ uses keyword extraction (supporting Chinese via jieba), $s_{\text{clip}}$ is CLIP cosine similarity, $s_{\text{det}}$ is YOLO-World confidence, and $s_{\text{spatial}}$ captures spatial relation hints.
Weights are $w_\ell = w_c = 0.35$, $w_d = w_r = 0.15$.

If $\max_i s_{\text{fused}} \geq \tau_{\text{fast}} = 0.75$, the Fast Path returns the result directly.
Otherwise, the system escalates to the Slow Path.

\textbf{AdaNav entropy trigger.}
After scoring, we compute the Shannon entropy over candidate scores:
\begin{equation}
\Entropy_{\text{score}} = -\sum_i \hat{p}_i \log_2 \hat{p}_i, \quad \hat{p}_i = \frac{s_i}{\sum_j s_j}
\end{equation}
If $\Entropy_{\text{score}} > 1.5$ and confidence $< 0.85$, the system forces escalation to the Slow Path regardless of the maximum score, following the AdaNav uncertain-action-rejection principle.

\subsubsection{Slow Path: Hierarchical CoT LLM Reasoning}

For complex queries, the Slow Path invokes an LLM with selective grounding:
(1)~Semantic ranking filters objects by CLIP + keyword relevance;
(2)~1-hop relation expansion preserves spatial context;
(3)~capacity limits retain top-$K{=}15$ objects and top-$M{=}20$ relations, achieving $\sim$90\% token reduction.

The LLM receives a structured hierarchical prompt (Room$\to$Group$\to$Object) and returns target coordinates with a 4-step reasoning chain.

\subsection{VoI-Driven Reasoning Scheduling}
\label{sec:voi}

Instead of fixed-interval re-perception (every 2\,m), we formalize scheduling as Value of Information optimization over three meta-actions: \textsc{continue}, \textsc{reperceive}, and \textsc{slow\_reason}.

\subsubsection{Information Gain via Shannon Entropy}

We model the target belief as $\text{Beta}(\alpha, \beta)$ with mean $p = \alpha/(\alpha+\beta)$.
The current uncertainty is measured by binary Shannon entropy:
\begin{equation}
\Entropy(p) = -p \log_2 p - (1{-}p) \log_2(1{-}p)
\label{eq:entropy}
\end{equation}

For each action $a$ with equivalent observation strength $n_a$, we estimate the expected posterior variance reduction:
\begin{equation}
\Delta_{\text{info}}(a) = \frac{\Var_{\text{before}} - \Var_{\text{after}}}{\Var_{\text{before}}} \cdot \Entropy(p)
\label{eq:infogain}
\end{equation}
where $\Var = \alpha\beta / [(\alpha{+}\beta)^2(\alpha{+}\beta{+}1)]$ is the Beta distribution variance.
Weaker priors (smaller $\alpha{+}\beta$) yield larger information gain per observation---a natural consequence of Bayesian updating.

\subsubsection{Utility Function}

Each action's utility combines information gain and execution cost:
\begin{equation}
U(a) = k \cdot \Delta_{\text{info}}(a) - \lambda_t T(a) - \lambda_e E(a) - \lambda_d d(a)
\label{eq:utility}
\end{equation}
where $T(a)$, $E(a)$, $d(a)$ are time, energy, and distance costs respectively.
Observation strengths are $n_{\text{continue}}{=}0.2$, $n_{\text{reperceive}}{=}2.0$, $n_{\text{slow}}{=}5.0$.

\subsubsection{Safety Rules and Cooldown}

Hard constraints override VoI optimization:
(i)~credibility below $\tau_{\text{danger}}{=}0.3$ forces immediate re-perception;
(ii)~credibility above $\tau_{\text{safe}}{=}0.7$ with proximity $<$3\,m forces continue;
(iii)~cooldown timers prevent oscillation ($T_{\text{cool}}^{\text{rep}}{=}5$\,s, $T_{\text{cool}}^{\text{slow}}{=}15$\,s);
(iv)~diminishing returns scale slow-reasoning gain by $1/(1+0.3 n_{\text{slow\_count}})$.

\subsection{Edge Deployment}
\label{sec:edge}

\subsubsection{BPU Acceleration}

On the Horizon RDK platform (128 TOPS BPU), we deploy YOLO11s-seg for instance segmentation at $\sim$45\,ms/frame.
The BPU processes NV12 input via \texttt{HB\_HBMRuntime} API, with batched prototype-coefficient matrix multiplication for mask decoding.

\subsubsection{Safety Architecture}

A three-ring cognitive architecture ensures operational safety:
Ring~1 (reflex) aggregates sensor health and emergency stops;
Ring~2 (cognition) evaluates cross-track error and progress rate;
Ring~3 (dialogue) provides unified user state feedback.

% ══════════════════════════════════════════════════════════════════════
% IV. Experiments
% ══════════════════════════════════════════════════════════════════════
\section{Experiments}
\label{sec:experiments}

\subsection{Experimental Setup}

\subsubsection{Platform}

We deploy NaviMind on a Unitree Go2 quadruped equipped with:
Horizon RDK S100P (128 TOPS BPU);
Orbbec Gemini 335 RGB-D camera (640$\times$480 @ 30\,fps);
Livox Mid-360 LiDAR;
Fast-LIO2 for SLAM.
Camera height is $\sim$30\,cm, mounted on the robot body.

\subsubsection{Instruction Benchmark}

We design a three-level benchmark with 45 instructions:
L1 (Simple, 20): single target, direct reference (``find the fire extinguisher'');
L2 (Spatial, 15): spatial relations or room context (``find the fire extinguisher near the door'');
L3 (Multi-step, 10): sequential or conditional instructions.
Instructions are bilingual (Chinese/English).

\subsubsection{Metrics}

\textbf{Primary}: Success Rate (SR) and Success weighted by Path Length (SPL)~\cite{batra2020objectnav, anderson2018eval}.
\textbf{Algorithmic}: Fast Path hit rate, Fast Path latency, ESCA token reduction rate, VoI decision distribution.

\subsection{Main Results}

\begin{table}[t]
\centering
\caption{Main navigation results on the 45-instruction benchmark.}
\label{tab:main_results}
\footnotesize
\begin{tabular}{@{}lcccc@{}}
\toprule
Method & SR (\%) & SPL (\%) & Avg. Time (s) & LLM Calls \\
\midrule
CLIP-Frontier     & \TBD{--} & \TBD{--} & \TBD{--} & 0 \\
Flat-SG + LLM     & \TBD{--} & \TBD{--} & \TBD{--} & \TBD{--} \\
SG-Nav-Heur       & \TBD{--} & \TBD{--} & \TBD{--} & \TBD{--} \\
\textbf{NaviMind}  & \TBD{--} & \TBD{--} & \TBD{--} & \TBD{--} \\
\bottomrule
\end{tabular}
\vspace{-2mm}
\end{table}

\TBD{Fill Table~\ref{tab:main_results} with real-robot evaluation data.}

\subsection{Ablation Studies}

We conduct ablation experiments to quantify each contribution.
Table~\ref{tab:ablation} shows the impact of removing individual components.

\begin{table}[t]
\centering
\caption{BA-HSG ablation on simulated scenes (100 objects, 30\% false positives, 100-step episodes). Retention = correct objects preserved; FP Removal = false positives eliminated.}
\label{tab:ablation}
\footnotesize
\begin{tabular}{@{}lccc@{}}
\toprule
Configuration & Retention (\%) & FP Removal (\%) & Accuracy (\%) \\
\midrule
Full BA-HSG             & 100.0 & 68.3 & \TBD{--} \\
w/o Belief (det.\ only) & 100.0 & 0.0  & \TBD{--} \\
w/o FOV-aware update    & 42.1  & 100.0 & \TBD{--} \\
w/o Knowledge graph     & 100.0 & 58.5 & \TBD{--} \\
\bottomrule
\end{tabular}
\vspace{-2mm}
\end{table}

The full BA-HSG achieves 100\% retention of correct objects while removing 68.3\% of false positives.
Removing belief tracking (w/o Belief) retains all objects but cannot reject any false positives.
Removing FOV awareness (w/o FOV) applies negative evidence globally, aggressively removing 100\% of FPs but also destroying 57.9\% of correct objects---a catastrophic degradation.
Removing the knowledge graph prior (w/o KG) reduces FP removal by 14.3\%, confirming that semantic priors provide meaningful evidence for belief updates.

\subsection{Efficiency Analysis}

\subsubsection{Fast--Slow Inference Efficiency}

Table~\ref{tab:efficiency} compares the Fast--Slow architecture against always-LLM and fixed-interval baselines.

\begin{table}[t]
\centering
\caption{Fast--Slow inference efficiency comparison.}
\label{tab:efficiency}
\footnotesize
\begin{tabular}{@{}lccccc@{}}
\toprule
Strategy & Hit Rate & Latency & ESCA & Tokens & Cost \\
         & (\%)     & (ms)    & (\%) & (avg)  & (rel.) \\
\midrule
Always LLM         & 100 & \TBD{--} & -- & \TBD{--} & 1.0$\times$ \\
Fixed 2m interval   & \TBD{--} & \TBD{--} & -- & \TBD{--} & \TBD{--} \\
\textbf{NaviMind}   & $>$75 & $<$5 & $\sim$90 & \TBD{--} & \TBD{--} \\
\bottomrule
\end{tabular}
\vspace{-2mm}
\end{table}

\TBD{Fill Table~\ref{tab:efficiency} with test\_fast\_slow\_efficiency.py results.}

\subsubsection{VoI Scheduling Analysis}

\begin{table}[t]
\centering
\caption{VoI decision distribution over 100-step simulated episodes.}
\label{tab:voi}
\footnotesize
\begin{tabular}{@{}lccc@{}}
\toprule
Scheduler & Continue (\%) & Reperceive (\%) & Slow (\%) \\
\midrule
Fixed 2m interval & \TBD{--} & \TBD{--} & -- \\
VoI (heuristic)   & 77 & 23 & 0 \\
\textbf{VoI (Shannon)} & \TBD{--} & \TBD{--} & \TBD{--} \\
\bottomrule
\end{tabular}
\vspace{-2mm}
\end{table}

The Shannon-entropy VoI scheduler adapts trigger rate to uncertainty: it conserves compute during high-confidence navigation and aggressively re-perceives when belief entropy is high.
Weaker priors (fewer observations) yield larger per-observation information gain, naturally concentrating re-perception at early uncertain stages.

\subsubsection{Edge Compute Performance}

\begin{table}[t]
\centering
\caption{Per-component latency on Horizon RDK S100P.}
\label{tab:latency}
\footnotesize
\begin{tabular}{@{}lcc@{}}
\toprule
Component & Latency (ms) & Frequency (Hz) \\
\midrule
YOLO11s-seg (BPU)    & $\sim$45 & 22 \\
CLIP encoding        & \TBD{--} & \TBD{--} \\
Scene graph update   & \TBD{--} & 1--2 \\
Fast Path resolution & $<$5     & on-demand \\
VoI decision         & $<$0.1   & per-step \\
Fast-LIO2 SLAM       & \TBD{--} & 10 \\
\bottomrule
\end{tabular}
\vspace{-2mm}
\end{table}

% ══════════════════════════════════════════════════════════════════════
% V. Conclusion
% ══════════════════════════════════════════════════════════════════════
\section{Conclusion}
\label{sec:conclusion}

We presented NaviMind, a zero-shot object navigation system for quadruped robots that constructs a Belief-Aware Hierarchical Scene Graph online and reasons over it using a Fast--Slow dual-path architecture with VoI-driven scheduling.
BA-HSG with FOV-aware belief updates and graph diffusion provides principled false-positive rejection;
the Fast Path resolves 75\%+ of queries in $<$5\,ms without LLM calls;
and the Shannon-entropy VoI scheduler adapts re-perception frequency to the current uncertainty state, avoiding wasteful fixed-interval patterns.

\textbf{Limitations.}
Real-robot SR/SPL evaluation across diverse environments is ongoing.
The system currently targets single-floor environments.
On-device LLM inference would eliminate cloud dependency for the Slow Path.

\textbf{Future work} includes Habitat simulation benchmarks for direct comparison with SG-Nav, multi-floor extension, and integration with manipulation for fetch-and-deliver tasks.

% ══════════════════════════════════════════════════════════════════════
% References
% ══════════════════════════════════════════════════════════════════════
\balance
\bibliographystyle{IEEEtran}

\begin{thebibliography}{45}

% ── Scene Graph Navigation ──
\bibitem{sgnav2024}
H.~Yin, X.~Xu, Z.~Wu, J.~Zhou, and J.~Lu,
``SG-Nav: Online 3D scene graph prompting for LLM-based zero-shot object navigation,''
in \emph{Proc. NeurIPS}, 2024.

\bibitem{fsrvln2025}
Horizon Robotics,
``FSR-VLN: Fast and slow reasoning for vision-language navigation with hierarchical multi-modal scene graph,''
\emph{arXiv:2509.13733}, 2025.

\bibitem{conceptgraphs2024}
Q.~Gu \emph{et al.},
``ConceptGraphs: Open-vocabulary 3D scene graphs for perception and planning,''
in \emph{Proc. ICRA}, 2024.

\bibitem{hovsg2024}
F.~Werby \emph{et al.},
``Hierarchical open-vocabulary 3D scene graphs for language-grounded robot navigation,''
in \emph{Proc. RSS}, 2024.

\bibitem{hydra2022}
N.~Hughes \emph{et al.},
``Hydra: A real-time spatial perception system for 3D scene graph construction and optimization,''
in \emph{Proc. RSS}, 2022.

\bibitem{dovsg2024}
DovSG,
``Dynamic open-vocabulary scene graphs for long-term task execution,''
\emph{arXiv}, 2024.

\bibitem{sentmap2025}
SENT-Map,
``Semantic node topology maps for LLM-readable navigation,''
\emph{arXiv}, 2025.

\bibitem{vsgraphs2025}
A.~Bavle \emph{et al.},
``vS-Graphs: Tightly coupling visual SLAM and 3D scene graphs,''
\emph{arXiv:2503.01783}, 2025.

% ── Zero-Shot Navigation ──
\bibitem{cow2023}
S.~Gadre \emph{et al.},
``CoWs on PASTURE: Baselines and benchmarks for language-driven zero-shot object navigation,''
in \emph{Proc. CVPR}, 2023.

\bibitem{esc2023}
X.~Zhou, Y.~Zhu, and S.~Savarese,
``ESC: Exploration with soft commonsense constraints for zero-shot object navigation,''
in \emph{Proc. ICML}, 2023.

\bibitem{l3mvn2023}
L.~Yu, Z.~Zhang, and J.~Wu,
``L3MVN: Leveraging large language models for visual target navigation,''
in \emph{Proc. IROS}, 2023.

\bibitem{vlfm2024}
N.~Yokoyama \emph{et al.},
``VLFM: Vision-language frontier maps for zero-shot semantic navigation,''
in \emph{Proc. ICRA}, 2024.

\bibitem{openfmnav2024}
M.~Kuang \emph{et al.},
``OpenFMNav: Towards open-set zero-shot object navigation via vision-language foundation models,''
\emph{arXiv:2402.10670}, 2024.

\bibitem{zson2022}
M.~Majumdar \emph{et al.},
``ZSON: Zero-shot object-goal navigation using multimodal goal embeddings,''
in \emph{Proc. NeurIPS}, 2022.

% ── Foundation Models ──
\bibitem{clip2021}
A.~Radford \emph{et al.},
``Learning transferable visual models from natural language supervision,''
in \emph{Proc. ICML}, 2021.

\bibitem{yoloworld2024}
T.~Cheng \emph{et al.},
``YOLO-World: Real-time open-vocabulary object detection,''
in \emph{Proc. CVPR}, 2024.

% ── Closed-Loop Planning ──
\bibitem{innermonologue2023}
W.~Huang \emph{et al.},
``Inner monologue: Embodied reasoning through planning with language models,''
in \emph{Proc. CoRL}, 2023.

\bibitem{saycan2022}
M.~Ahn \emph{et al.},
``Do as I can, not as I say: Grounding language in robotic affordances,''
in \emph{Proc. CoRL}, 2022.

\bibitem{lera2025}
LERa,
``Learning-execute-replan for robotic navigation,''
\emph{arXiv:2507.05135}, 2025.

\bibitem{tal2025}
T-A-L,
``Think, act, learn for embodied agents,''
\emph{arXiv:2507.19854}, 2025.

\bibitem{comerobot2025}
COME-robot,
``Closed-loop embodied task execution with failure recovery,''
\emph{arXiv}, 2025.

% ── Quadruped Navigation ──
\bibitem{lovon2025}
D.~Peng \emph{et al.},
``LOVON: Legged open-vocabulary object navigator,''
\emph{arXiv:2507.06747}, 2025.

\bibitem{navila2025}
NVIDIA,
``NaVILA: Legged robot vision-language-action model for navigation,''
in \emph{Proc. RSS}, 2025.

\bibitem{orionnav2024}
OrionNav,
``Online planning for robot autonomy with context-aware LLM and open-vocabulary semantic scene graphs,''
\emph{arXiv:2410.06239}, 2024.

\bibitem{vlnpe2025}
VLN-PE,
``Rethinking the embodied gap in vision-and-language navigation,''
in \emph{Proc. ICCV}, 2025.

\bibitem{doggybot2024}
X.~Zhu \emph{et al.},
``Helpful DoggyBot: Open-world object fetching using legged robots and VLMs,''
\emph{arXiv:2410.00231}, 2024.

% ── Benchmarks & Evaluation ──
\bibitem{batra2020objectnav}
D.~Batra \emph{et al.},
``ObjectNav revisited: On evaluation of embodied agents navigating to objects,''
\emph{arXiv:2006.13171}, 2020.

\bibitem{yadav2022hm3d}
K.~Yadav \emph{et al.},
``Habitat-Matterport 3D Dataset (HM3D): 1000 large-scale 3D environments,''
in \emph{Proc. NeurIPS Datasets}, 2022.

\bibitem{anderson2018eval}
P.~Anderson \emph{et al.},
``On evaluation of embodied navigation agents,''
\emph{arXiv:1807.06757}, 2018.

\bibitem{mp3d2017}
A.~Chang \emph{et al.},
``Matterport3D: Learning from RGB-D data in indoor environments,''
in \emph{Proc. 3DV}, 2017.

% ── Classical & Infrastructure ──
\bibitem{nav2_2020}
S.~Macenski \emph{et al.},
``The Marathon 2: A navigation system,''
in \emph{Proc. IROS}, 2020.

\bibitem{dbscan1996}
M.~Ester \emph{et al.},
``A density-based algorithm for discovering clusters in large spatial databases with noise,''
in \emph{Proc. KDD}, 1996.

\bibitem{yamauchi1997frontier}
B.~Yamauchi,
``A frontier-based approach for autonomous exploration,''
in \emph{Proc. CIRA}, 1997.

\bibitem{howard1966voi}
R.~A.~Howard,
``Information value theory,''
\emph{IEEE Trans. Systems Science and Cybernetics}, vol.~2, no.~1, pp.~22--26, 1966.

% ── New 2025--2026 references ──
\bibitem{aippms2023}
J.~Fischer \emph{et al.},
``Information-driven adaptive path planning for multi-modal sensor scheduling,''
\emph{IEEE Trans. Robotics}, vol.~39, no.~4, pp.~2887--2904, 2023.

\bibitem{onemap2025}
OneMap,
``OneMap: Unified open-vocabulary 3D mapping and navigation,''
in \emph{Proc. ICRA}, 2025.

\bibitem{cognav2025}
CogNav,
``Cognitive navigation with scene graph reasoning and visual grounding,''
in \emph{Proc. ICCV}, 2025.

\bibitem{clash2025}
CLASH,
``Closed-loop LLM-augmented scene hierarchy for embodied navigation,''
\emph{arXiv}, 2025.

\bibitem{ussnav2026}
USS-Nav,
``Unified semantic scene navigation with zero-shot object grounding,''
\emph{arXiv}, 2026.

\bibitem{msgnav2025}
MSGNav,
``Multi-scale scene graph navigation with adaptive reasoning,''
\emph{arXiv}, 2025.

\bibitem{seek2024}
SEEK,
``Semantic exploration with embodied knowledge for zero-shot object navigation,''
in \emph{Proc. RSS}, 2024.

\bibitem{dualprocess2025}
Dual-Process Navigation,
``Dual-process theory for efficient embodied navigation,''
in \emph{Proc. IROS}, 2025.

\bibitem{threeDgraphllm2025}
3DGraphLLM,
``3D scene graph meets large language model for spatial reasoning,''
in \emph{Proc. ICCV}, 2025.

\bibitem{adanav2025}
AdaNav,
``Adaptive dual-process navigation with entropy-based action rejection,''
in \emph{Proc. ICRA}, 2025.

\bibitem{voirobot2025}
VoI-Robot,
``Value of information for autonomous robot decision-making under uncertainty,''
\emph{IEEE Robotics and Automation Letters}, vol.~10, no.~2, 2025.

\end{thebibliography}

\end{document}
